% Q08_gold_answers.tex — GRADER CALIBRATION FIXTURE (a correct answer set).
% A drop-in replacement for paper/Q08_answers.tex (spec §4.7, §6.12), written in a
% learner's words, paraphrased away from the reference solutions. Every claim was
% checked with the rubric's probes on the course pins. Reference-class (hook rule 1,
% DELTAS Δ125): it is a correct answer set, so open it only after Q08 is done.
%
% EXPECTED[B1][recomputation][4]: correct (Probe A's TermIs verifies on the final term; Probe A''s Steps passes on the three steps; Probe B's MeaningIs verifies on min-form 0 <= s(0) && 0 <= s(1); Probe C's SwappedMeaningIs verifies on s(2) - s(0) <= s(1), and SwappedCounterexample passes on (2, -3, -9)).
% EXPECTED[B2][recomputation][4]: correct (Probe D verifies with the answer's maintenance formula; Probe E's HaltState passes on x = -1, the answer's form "non-negative while the guard holds", and HaltTrace on 3, 1, -1, -3, -5; Probe F's runs match (c)).
% EXPECTED[B3][card-fidelity][4]: every field agrees with lines 584-610; the contract mapping is defensible (CF.g).
% Expected per-problem grades: B1 = 4, B2 = 4, B3 = 4. A finding <= 2 on any problem is
% a calibration failure (a hallucinated flaw); a 3 needs a named gap that the answer
% really has.
\documentclass[11pt]{article}
\input{preamble}

\title{Q08 --- Internals II: refinement, permissions, verified wp \\ \large Written answers}
\author{learner}
\date{}

\begin{document}
\maketitle

% ---------------------------------------------------------------------------
\section*{Problem B1 (hand-calculation)}\label{ans:B1}
\begin{answer}
\textbf{(a)} \dfy{Wp} on a \dfy{Seq} works back to front: first the last
assignment, then the \dfy{If} against what that gives. \\
Assignment \dfy{Assign(0, Minus(Var(2), Var(0)))} (\dfy{Wp} case
\dfy{Assign}, so \dfy{SubstB} on \dfy{Le}, so \dfy{SubstE} on each side:
\dfy{Var(0)} is the substituted variable and becomes the expression,
\dfy{Var(1)} is not and stays): \dfy{Le(Minus(Var(2), Var(0)), Var(1))}. Call
it $p$. \\
The \dfy{If} (\dfy{Wp} case \dfy{If}): \dfy{And(Imp(b, A), Imp(Not(b), E))}
where $b$ is \dfy{Le(Var(0), Var(1))} as it stands, $A$ is
\dfy{Wp(Assign(2, Var(1)), p)} and $E$ is \dfy{Wp(Assign(2, Var(0)), p)}. \\
$A$: \dfy{SubstB} on \dfy{Le}, \dfy{SubstE} through \dfy{Minus} into the
variables; only \dfy{Var(2)} changes, to \dfy{Var(1)}:
\dfy{Le(Minus(Var(1), Var(0)), Var(1))}. \\
$E$: likewise, \dfy{Var(2)} becomes \dfy{Var(0)}:
\dfy{Le(Minus(Var(0), Var(0)), Var(1))}. \\
So \dfy{Wp(C(), Q())} is
\begin{dafnycode}
And(Imp(Le(Var(0), Var(1)), Le(Minus(Var(1), Var(0)), Var(1))),
    Imp(Not(Le(Var(0), Var(1))), Le(Minus(Var(0), Var(0)), Var(1))))
\end{dafnycode}

\textbf{(b)} Reading it with \dfy{Holds} and \dfy{Eval}: if $s(0) \leq s(1)$
then $s(1) - s(0) \leq s(1)$, and if not then $s(0) - s(0) \leq s(1)$. \\
The first conclusion is $s(0) \geq 0$, the second $s(1) \geq 0$; so: when $x
\leq y$ we need $x \geq 0$ (and then $y \geq 0$ follows), when $x > y$ we need
$y \geq 0$ (and then $x \geq 0$ follows). Both cases come to
$0 \leq s(0) \wedge 0 \leq s(1)$: the smaller of $x$ and $y$ is non-negative.
\\
It never mentions $s(2)$: whatever $m$ starts as, the \dfy{If} overwrites it
before $x := m - x$ uses it.

\textbf{(c)} If \dfy{Seq} composed the other way round, the \dfy{If} would see
$q$ itself; assigning $m$ cannot affect $x \leq y$, so that stage means $x
\leq y$, and substituting $m - x$ for $x$ afterwards (into the condition as
well) leaves $s(2) - s(0) \leq s(1)$. \\
Start state: $x = 2$, $y = -3$, $m = -9$; $-9 - 2 = -11 \leq -3$, so the
swapped rule accepts it. \\
Run: $2 \leq -3$ is false, so $m$ becomes $2$; then $x$ becomes $2 - 2 = 0$;
we end with $x = 0$ and $y = -3$, and $0 \leq -3$ fails. \\
What should have warned me: a precondition that constrains $m$'s initial value
even though the program never reads it, and a guard rewritten as if it ran
after the assignment.
\end{answer}

% ---------------------------------------------------------------------------
\section*{Problem B2 (hand-calculation)}\label{ans:B2}
\begin{answer}
\textbf{(a)} Entry: after $x := x_0$ and $k := 0$, $I$ reads $0 + \lvert
x_0\rvert = \lvert x_0\rvert$: always true. \\
Maintenance: back through $k := k + 1$ the target becomes $k + 1 + \lvert
x\rvert = \lvert x_0\rvert$ and $0 \leq \lvert x\rvert < D$; back through the
conditional it splits: \\
\hspace*{2em} when $x > 0$: $k + 1 + \lvert x - 1\rvert = \lvert x_0\rvert$ and
$0 \leq \lvert x - 1\rvert < D$. Here $x - 1 \geq 0$, so $\lvert x - 1\rvert$
is one less than $\lvert x\rvert = D$: both follow from $I$. \\
\hspace*{2em} when $x \leq 0$: $k + 1 + \lvert x + 1\rvert = \lvert x_0\rvert$
and $0 \leq \lvert x + 1\rvert < D$. Since $x \neq 0$ here, $x$ is negative and
$x + 1 \leq 0$, so $\lvert x + 1\rvert$ is again $\lvert x\rvert - 1$: both
follow. \\
Use at exit: the guard is false, so $x = 0$ and $I$ says $k = \lvert
x_0\rvert$, the \dfy{ensures}. \\
The guard does work twice: in the $x \leq 0$ branch, to exclude $x = 0$ (there
one more step would give $\lvert x\rvert = 1$, wrong for both $I$ and the
decrease), and, negated, at exit.

\textbf{(b)} Soundness: induct on how many iterations have run: before the
first test $I$ holds by entry, and if $I$ holds at a test the loop passes,
maintenance gives it at the next test. So $I$ holds whenever the loop stops,
together with $\neg b$, and that gives the postcondition. Each iteration
lowers the integer measure by at least one, and it is never negative while
the loop runs, so the loop cannot run forever. \\
For \dfy{Halt}, entry, exit and the decrease are fine; the bound is what
fails. Writing it as ``the measure is non-negative while the guard holds'':
at $x = -1$ the guard holds and the measure is $-1$. \\
From $3$: $3, 1, -1, -3, -5, \ldots$, forever.

\textbf{(c)} \vrs{x - 1} is only done when $x > 0$, so it stays inside
\vrs{i64}; \vrs{x + 1} only when $x \leq 0$, so it stays inside too. For
\vrs{k + 1}: the invariant and the guard say $k = \lvert x_0\rvert - \lvert
x\rvert$ with $\lvert x\rvert \geq 1$, and $\lvert x_0\rvert$ is at most the
magnitude of \vrs{i64::MIN}, so $k + 1$ fits in \vrs{u64} with room to spare.
\\
With $k$ as \vrs{u32}: it fails at \vrs{k + 1}, since $x_0$ can be far larger
than \vrs{u32::MAX}, and the invariant then allows $k$ $=$ \vrs{u32::MAX}
while $x \neq 0$.
\end{answer}

% ---------------------------------------------------------------------------
\section*{Problem B3 (paper card)}\label{ans:B3}
\begin{answer}
\textbf{Source:} cheatsheet entry 02.1, lines 584--610 (the pinned range).

\textbf{Task shape:} One large proof: that CertiCoq's ANF transformation
(untyped de Bruijn source to a named ANF IR, both with fuelled big-step
semantics) preserves the meaning of terminating programs, written in Rocq by
the agent with the existing CPS proof as its template.

\textbf{Oracle:} The Rocq compiler accepting the development; the agent only
edited files and ran \texttt{rocqc}. That accepts whatever statements are in
the files, so a silently weakened statement passed it: an expert had to watch
the statements.

\textbf{What is trusted:} Rocq itself; the definitions and the theorem
statements (semantics with fuel, the logical relation); the expert steering
the run; and the two tactics a human repaired.

\textbf{Cheat surface:} Statement weakening (a fuel bound dropped, with a
trivial zero for the fuel obligation), deleting a proof so a file compiles,
and a false assumption that two result variables cannot be equal. \\
In my flagship contract: \texttt{specfreeze --against} catches the first; the
gate (zero holes, a full run) catches the second; the third only
\texttt{vaudit} catches if it was an \texttt{assume}-style axiom, and otherwise
only a reader or a witness, and the range does not say which it was.

\textbf{Metric:} Did the proof go through (yes, except divergence
preservation), plus size and effort against the human CPS proof: 7,783 vs
5,294 lines in all, 4,755 vs 2,243 lines for the main theorem, and the cost
in hours, tool calls, compiles and compactions.

\textbf{Budget:} Nothing fixed beforehand; what it used: 114 prompts by the
expert, 5 days (90.2 h wall-clock, 62.6 h active, about 96 h by the text),
6,538 tool calls, 872 compiles, 101 compactions; the expert's own time is not
measured.

\textbf{Variance:} Not addressed: one run, one compiler pass, one template,
one expert (n = 1); the line counts in the text and the table disagree.

\textbf{Artifact:} An agent-written Rocq correctness proof for CertiCoq's ANF
pass, reported in PACMPL 2026 (arXiv 2602.20082); the model was Claude Opus 4.6
under Claude Code. The range is silent on a released artifact or transcripts.
\end{answer}

\end{document}
